随机游走须先写清吸收、反射和重启规则；方向有记忆或边界随历史移动时，应把方向或回撤差距纳入状态。命中概率满足调和方程，平均时间满足带常数项的首步方程。对成功路径作条件化时，要按后续成功概率重新加权，不能只删去失败终点。

鞅提供另一类边值工具，但无界停时不能直接代入期望恒等式。可先停在 $\tau\wedge n$，再凭一致可积性或可积支配取极限；仅有几乎必然停止并不够。跳跃可能越过边界，终值也须按真实越界位置处理。
\par\smallskip
